\begin{textsample}{POS Dim 4 – rn – Score 12.00 – rn/rn\_em\_5.txt}  \label{ex:f4_pos_084}
1.5 SUJEITOS DO ENSINO MÉDIO POTIGUAR Identificar sujeitos, segundo alguns autores como Karl Mannhein ( 1976 ), parte de um princípio do qual os indivíduos formam suas personalidades intimamente relacionadas onde se realizam suas interações sociais.
Dessa forma, traçar um perfil das juventudes potiguares, remete -nos a aspectos histórico-culturais e socioeconômicos que nos auxiliam na compreensão dos cômputos educacionais que, rotineiramente, são atribuídos exclusivamente à falta de políticas públicas direcionadas a este segmento, cabendo aqui neste Referencial Curricular o dever de contribuirmos com estratégias didáticas que minimizem as discrepâncias destes impactos.
O \textbf{Estado} do Rio Grande do Norte é, conforme o IBGE, 2018, composto por aproximadamente 3 ( três ) \textbf{milhões} e meio de habitantes.
Essa \textbf{população} ocupa os 167 ( cento e sessenta e sete ) municípios, com uma ligeira predominância de representantes do sexo feminino. \textbf{Pardos} 52,25 \% e \textbf{Brancos} 40,84 \% compõem a maioria dos grupos \textbf{raciais}, seguidos por pretos 5,23 \% e \textbf{indígenas} 0,09 \%.
Tal distribuição é resultante da descendência entrelaçada de três povos : os negros, \textbf{indígenas} e portugueses, respectivamente, povos escravizados, nativos e colonizadores.
Há registro de 33 ( trinta e três ) comunidades quilombolas em 26 ( vinte e seis ) cidades desta unidade federativa e 11 ( onze ) comunidades \textbf{indígenas} reconhecidas, distribuídas entre os municípios.
Destacamos a vulnerabilidade social a qual estão submetidos os dois últimos grupos, dificultando o desenvolvimento pleno de atividades educativas.
Ainda nesse contexto, constatamos uma distribuição geográfica agrupada nos centros urbanos ( 82 \% ).
Contudo, essa disposição está concentrada a um número reduzido de cidades, sendo a maioria delas composta por uma divisão mais equitativa da \textbf{população} urbe-campestre, sobretudo no interior, com uma parcela significativa delas com prevalência de \textbf{população} campesina, o que pode ser considerado um obstáculo para os nossos jovens quanto ao acesso aos meios de comunicação, conexão global, tecnologia, transportes, energia elétrica e saneamento básico.
Somando -se a este fato, devemos considerar também fatores climáticos \textbf{marcantes} em nosso solo, no qual encontra -se localizado 90 \% de sua totalidade na região do Polígono das Secas, com abrangência do clima tropical \textbf{semiárido}, perpassando praticamente todo o interior do \textbf{estado} e \textbf{litoral} \textbf{norte}.
Vislumbramos, nesse cenário, um outro entrave aos nossos estudantes, pois muitos têm que se desdobrarem em duplas jornadas de trabalho e estudos, muitas vezes ausentando -se das instituições escolares por períodos sazonais que coincidem com colheitas ou ciclos econômicos microrregionais periódicos.
Este panorama de carência econômica familiar atinge 66 \% dos alunos matriculados na rede estadual de ensino que possuem renda média familiar de até 2 ( dois ) salários mínimos, evidenciando a necessidade do trabalho precoce.
O \textbf{Estado} do Rio Grande do Norte é dotado de um grande número de jovens, fruto de altas taxas de natalidade e fertilidade históricas, embora venha declinando nas últimas décadas.
A faixa etária de 15 ( quinze ) a 29 ( vinte e nove ) anos representa 1⁄4 da \textbf{população}, dos quais 15 \% estão matriculados na rede estadual de ensino.
Gráfico 2 – Pirâmide etária Fonte : IBGE ( 2010 ).
Além do trabalho prematuro de ambos os sexos já elencados, chamamos atenção a grande quantidade de estudantes do sexo feminino fora da escola e sem exercer atividade remunerada.
Dados apontam que muitos adolescentes abandonam a escola em virtude da gravidez precoce e afazeres familiares e domésticos.
Um cenário que abrange todo o país, mas que se agrava em nossa região devido a estereótipos arraigados culturalmente.
Uma vez traçado o perfil dos sujeitos das nossas juventudes, ponderando seus aspectos históricos, culturais e socioeconômicos, propiciamos a adoção de diretrizes que contemplem seus anseios e potencialidades, respeitando suas adversidades, propondo ações e políticas que assegurem os direitos constitucionais educacionais desses jovens, preparando -os para os desafios do mundo pessoal e profissional diante do contemporâneo e cada vez mais exigente processo de transformação global e informacional, visando a formação de sujeitos ativos no exercício de sua cidadania.

% matched lemmas: branco, estado, indígena, litoral, marcante, milhão, norte, pardo, população, racial, semiárido
\end{textsample}
